\documentclass[11pt,letterpaper,reqno]{amsart}
\usepackage[T1]{fontenc}
\usepackage{lmodern}
\usepackage[margin=1in]{geometry}
\usepackage{amsmath,amssymb,amsthm,mathtools,mathrsfs}
\usepackage{booktabs,microtype}
\usepackage[colorlinks=true,linkcolor=red,citecolor=red,urlcolor=blue,
  pdfauthor={Simon Becker and Benjamin Lawton},pdftitle={Counterexamples to the ordinary Fried conjecture near hyperbolic metrics},
  pdfsubject={Counterexamples to the ordinary Fried conjecture},
  pdfkeywords={Fried conjecture, Ruelle zeta function, dynamical torsion}]{hyperref}
\newtheorem{theorem}{Theorem}[section]
\newtheorem{lemma}[theorem]{Lemma}
\newtheorem{proposition}[theorem]{Proposition}
\theoremstyle{remark}
\newtheorem{remark}[theorem]{Remark}
\newcommand{\LL}{\mathcal L}
\newcommand{\Id}{\operatorname{Id}}
\newcommand{\Res}{\operatorname{Res}}
\newcommand{\tr}{\operatorname{tr}}
\newcommand{\str}{\operatorname{str}}
\newcommand{\Hol}{\operatorname{Hol}}
\newcommand{\Det}{\operatorname{Det}}
\newcommand{\ran}{\operatorname{ran}}
\newcommand{\vol}{\operatorname{vol}}
\newcommand{\cH}{\mathcal H}
\newcommand{\cU}{\mathcal U}
\newcommand{\cV}{\mathcal V}
\newcommand{\cC}{\mathcal C}
\title[Counterexamples to the ordinary Fried conjecture]
{Counterexamples to the ordinary Fried conjecture\newline near hyperbolic metrics}
\author{Simon Becker}
\email{simon.becker@unibocconi.it}
\address{Bocconi University, Milan, Italy}
\author{Benjamin Lawton}
\address{Paraform, San Francisco, USA}
\date{12 October 2026}
\keywords{Fried conjecture, Ruelle zeta function, dynamical torsion, contact Anosov flow}

\begin{document}
\begin{abstract}
We report on an AI-assisted construction of counterexamples to the ordinary Fried conjecture
with real analytic metrics arbitrarily close to a hyperbolic
metric in dimension three. For acyclic unitary
line characters tending to the trivial character, the modulus of
the Ruelle zeta value at zero divided by analytic torsion squared
tends to zero. The proof
uses a resonance crossing, holomorphic dependence on the character, and
torsion variation on a fixed spectral cut.
\end{abstract}
\maketitle
\section{The statement}
In this article, we give a counterexample to a torsion identity conjectured by Fried for real analytic negatively
curved manifolds~\cite[pp.~180--181]{Fried95}.
We consider a closed oriented three-manifold $(\Sigma,g)$ with a real
analytic negatively curved metric and a unitary character
$\chi:\pi_1(\Sigma)\to U(1)$, where $U(1)$ is the unit circle.
Write $H^k(\Sigma;\chi)$ for cohomology with coefficients in its flat
line bundle, and $H^\bullet$ for all degrees. Acyclicity means
\[
 H^k(\Sigma;\chi)=0\qquad\text{for every }k.
\]
We write $\mathcal G_g^\#$ for the
oriented periodic orbits of the geodesic flow that are not repetitions
of shorter periodic orbits, and
$\ell_g(\gamma)$ for the length of an orbit $\gamma$.
For such an orbit, $\chi(\gamma)$ is the holonomy along its projection
to $\Sigma$.
For a complex variable $s$, both orientations of every closed geodesic
occur in the product
\begin{equation}\label{Euler}
 Z_{g,\chi}(s)=\prod_{\gamma\in\mathcal G_g^\#}
       \bigl(1-\chi(\gamma)e^{-s\ell_g(\gamma)}\bigr),
 \qquad \operatorname{Re}s\gg1.
\end{equation}
The product defines the Ruelle zeta function for sufficiently large
$\operatorname{Re}s$. We use its meromorphic continuation to define its
value near zero; see~\cite{DGRS}.
We write $T_\chi(\Sigma)$ for Ray--Singer analytic torsion~\cite{RS}.
For acyclic unitary coefficients it is independent of the metric.
In dimension three, Fried's formulation for acyclic unitary line
characters is
\begin{equation}\label{Friedidentity}
 |Z_{g,\chi}(0)|=T_\chi(\Sigma)^2.
\end{equation}
Zworski's Conjecture~9~\cite{Zw} states this identity for smooth
negatively curved metrics. We use $i$ for the imaginary unit and
$\lambda$ for a complex spectral variable; his convention is
$\zeta_\chi(\lambda)=Z_{g,\chi}(-i\lambda)$.
See also \cite[Conjecture~21, (35)]{HKS} and
\cite[Introduction, (1.2)]{DGRS} for versions of this conjecture.
Our construction stays in the real analytic class.

Write $\mathbf1$ for the trivial character, our main result is then: 
\begin{theorem}\label{mainresult}
There exist a closed oriented hyperbolic three-manifold $(\Sigma,g_0)$,
real analytic negatively curved metrics $g_j\to g_0$ in $C^\infty$,
and acyclic unitary
line characters $\chi_j\to\mathbf1$ such that $Z_{g_j,\chi_j}$ is
holomorphic at zero and
\begin{equation}\label{mainratio}
 \frac{|Z_{g_j,\chi_j}(0)|}{T_{\chi_j}(\Sigma)^2}\longrightarrow0.
\end{equation}
Thus \eqref{Friedidentity} fails for all sufficiently large $j$.
\end{theorem}

Using the untwisted splitting in~\cite{CDDP}, we construct a twisted
resonance crossing. Holomorphic character dependence, resonant exactness, and fixed-cut torsion
variation control the value there; Fried's theorem~\cite{Fried} normalizes it.
For each sufficiently small real $\theta\ne0$, \eqref{Friedidentity}
holds at every other metric of our small family, as proved in
Section~\ref{sec:normalization}; the selected metrics are its unique crossings.
The acyclic neighborhoods in~\cite[Theorem~3]{DGRS} depend on the fixed
character; our characters tend to the nonacyclic trivial character.


\textbf{AI statement.}
S.B. reports obtaining the original argument from the single
prompt \emph{"Prove Conjecture 9"}, with reference \cite{Zw} attached, using
GPT-6.1 Sol with Ultra reasoning in Codex. The first version of this manuscript referenced the repository of B.L. which ultimately led to this collaboration. AI assistance was also used in the preparation and revision of this manuscript. The original
\href{https://github.com/SimonLukas2604/Fried-s-conjecture/blob/0dcb387a3451c6d4d173ccc38b7c8d18c21ac462/original-gpt-outputs/holomorphic_continuity_crossing.tex}{manuscript output}
and
\href{https://github.com/SimonLukas2604/Fried-s-conjecture/blob/0dcb387a3451c6d4d173ccc38b7c8d18c21ac462/original-gpt-outputs/specialization_nilpotent_residue_audit_20261011.tex}{specialization audit output}
are on github. A comprehensive version, together with Lean verification of the
conditional algebraic and limiting steps, is available on
\href{https://github.com/SimonLukas2604/Fried-s-conjecture/tree/0dcb387a3451c6d4d173ccc38b7c8d18c21ac462}{github}.
The geometric and spectral results used from the literature are not formalized.

B.L. reports substantial Claude and Claude Code assistance
with research, derivation, writing, checks, and Lean formalization. His
\href{https://github.com/blawton/fried-conj-lean/blob/610d659ea15af869553f6e3e79c432880d9b0ab0/docs/counterexample_intuitive_view.tex}{disclosure}
and \href{https://github.com/blawton/fried-conj-lean/blob/610d659ea15af869553f6e3e79c432880d9b0ab0/formalization.yaml}{metadata}
state the earlier developments. The Lean theorem is conditional, as the geometric hypotheses are not formalized.

Both authors directed the work, checked the sources and mathematical content, and take responsibility for
the mathematical content of this work. 


\noindent\textbf{Acknowledgment.} S.B. thanks Maciej Zworski for a
discussion about the Fried conjecture that led him to explore this problem. 


\section{The spectral cut and the crossing}\label{sec:spectral}
For each degree $k$, we write $b_k(\Sigma)$ for the dimension of
$H^k(\Sigma;\mathbb R)$, the cohomology with real coefficients.
We choose $(\Sigma,g_0)$ with first Betti number $b_1(\Sigma)=1$ and a
nonzero real harmonic one-form $\omega$. In Appendix~\ref{app:inputs},
we construct such a base and a real analytic function $b$ with the
required conformal splitting. We use $d$ for the exterior differential
and $\wedge$ for the exterior product of forms. For a complex
character parameter $z$ and a small real metric parameter
$\tau\geq0$, we set
\[
 d_z=d+iz\omega\wedge,\qquad
 \chi_z(\gamma)=\exp\left(-iz\int_\gamma\omega\right),\qquad
 g_\tau=e^{-2\tau b}g_0.
\]
For small $\tau$, these metrics remain negatively curved, so their
geodesic flows are Anosov.
For every sufficiently small real $z=\theta\ne0$, $\chi_\theta$
is acyclic and unitary, as proved in Appendix~\ref{app:inputs}.
We work on the unit tangent bundle $M=S_{g_0}\Sigma$ of $g_0$.
For $(x,v)\in M$, we define $F_\tau(x,v)=(x,e^{\tau b(x)}v)$.
This map identifies $M$ with the unit tangent bundle of $g_\tau$.
We write $\pi:M\to\Sigma$ for the projection and $\pi^*$ for pullback.
On $M$, $d_z$ denotes $d+iz\pi^*\omega\wedge$.
We denote the pulled-back geodesic contact form and Reeb field by
$\alpha_\tau$ and $X_\tau$. For a vector field $Y$, we write $\iota_Y$ for
contraction by $Y$. In particular,
$\alpha_\tau(X_\tau)=1$ and $\iota_{X_\tau}d\alpha_\tau=0$.
We set
\[
 L_{z,\tau}=d_z\iota_{X_\tau}+\iota_{X_\tau}d_z,
 \qquad B_{z,\tau}=-L_{z,\tau}.
\]
We suppress parameters on operators when clear. We write $\Id$ for the
identity, $\tr$ for the trace, and $\ran$ for the range of an operator.
Resonances in the spectral
coordinate $s$ are poles of the continued resolvent. On the finite
spaces used below they are eigenvalues of $B$, since
\[
 (L+s\Id)^{-1}=(s\Id-B)^{-1}.
\]
The convention in~\cite{CDDP} is $P=-iL=iB$.

We choose a small radius $0<\lambda_{\rm cut}<1$ and fix the positively
oriented circle $\gamma_{\rm cut}=\{|s|=\lambda_{\rm cut}\}$.
We call this fixed circle our spectral cut.
It encloses only the zero resonance at $(z,\tau)=(0,0)$ and has no
boundary resonance in any degree. For $0\leq k\leq4$, we write
$\cU^k_{z,\tau}$ for the
generalized resonance space inside this circle on horizontal
degree-$k$ forms $u$, meaning $\iota_{X_\tau}u=0$.
A generalized space includes the Jordan vectors as well as the eigenvectors:
on the part for an eigenvalue $\lambda$, some power of
$B-\lambda\Id$ vanishes.
Every degree-$k$ form has a unique decomposition into a horizontal
part $u$ and a part containing $\alpha_\tau$:
\[
 u+\alpha_\tau\wedge v,\qquad
 \iota_{X_\tau}u=\iota_{X_\tau}v=0.
\]
Here $v$ has degree $k-1$. We use ``full forms'' when we include both parts.
The horizontal ranks in degrees $0,\ldots,4$, and the full-form
ranks in degrees $0,\ldots,5$, are
\begin{equation}\label{ranks}
 (1,2,4,2,1),\qquad (1,3,6,6,3,1).
\end{equation}
We use the fixed cut to follow a finite-dimensional cluster as the
parameters vary. These ranks count algebraic multiplicities, including
possible Jordan blocks. The initial ranks are verified in
Appendix~\ref{app:inputs}; Lemma~\ref{holom} proves that they persist.
Individual
resonances may cross zero without changing the dimension of the cluster.
We define
\[
 A_k=B|_{\cU^k},\qquad c=\tr A_0,\qquad t=\tr A_1-c,
 \qquad P_k(s,z,\tau)=\det(s\Id-A_k).
\]
Since $\cU^0$ has dimension one, $A_0=c\Id$.
For a real character parameter $\theta$, we write $d_\theta^*$ for the
$L^2(g_0)$ adjoint of $d_\theta$ and
\[
 \Delta_{\mathrm H,\theta}^{(k)}=d_\theta^*d_\theta+d_\theta d_\theta^*
\]
for the twisted Hodge Laplacian on degree-$k$ forms on $\Sigma$.
We denote its small scalar eigenvalue by $\mu_0(\theta)$ and set
\[
 c_0(\theta)=-1+\sqrt{1-\mu_0(\theta)},
\]
with the square root equal to one at the origin. The $L^2$ norm uses $g_0$, and
$\vol(\Sigma)$ denotes its volume. We set
$\kappa=\|\omega\|_{L^2(\Sigma)}^2/\vol(\Sigma)>0$.
For small real $\theta\ne0$, the calculations in Appendix~\ref{app:inputs},
in particular \eqref{app:scalarquadratic} and
\eqref{app:clusterpolynomials}, give
\begin{align}\label{hyper}
 c(\theta,0)=t(\theta,0)&=c_0(\theta)<0,
 &P_2(s,\theta,0)&=(s-c_0)^3(s+c_0),\\
 \mu_0(\theta)&=\kappa\theta^2+O(\theta^4),
 &c_0(\theta)&=-\tfrac\kappa2\theta^2+O(\theta^4).\label{quadratic}
\end{align}
The cut derivative computed in Appendix~\ref{app:inputs} gives
$r=\tr B'_{\rm cut}>0$ and
\begin{equation}\label{taxis}
 c(0,\tau)=0,\qquad t(0,\tau)=r\tau+o(\tau).
\end{equation}

\begin{lemma}\label{holom}
For some fixed $\rho,\varepsilon>0$, the cut matrices and their
characteristic coefficients are holomorphic for $|z|<\rho$.
They are $C^1$ in $\tau\in[0,\varepsilon]$. Their metric derivatives
are holomorphic in $z$ on the same disk and jointly continuous,
uniformly on each smaller closed disk.
The coefficients are even in $z$ and real for real $z$.
\end{lemma}
\begin{proof}
We take the common spaces $\cH_0,\ldots,\cH_4$ from
\cite[Appendix B.2--B.6]{CD}, on which the continued resolvent is bounded.
Their inclusions are continuous:
\[
 \cH_4\hookrightarrow\cH_3\hookrightarrow\cH_2
       \hookrightarrow\cH_1\hookrightarrow\cH_0.
\]
Moving from $\cH_{j+1}$ to $\cH_j$ loses one regularity level.
In each degree the untwisted resolvents
$R_{0,\tau}(s)=(L_{0,\tau}+s)^{-1}$ are uniformly bounded on the
contour at every level. The twist is a bounded multiplier:
\[
 L_{z,\tau}=L_{0,\tau}+zV_\tau,
 \qquad V_\tau=i\pi^*\omega(X_\tau)
 \quad\text{\cite[(8.9)]{CD}}.
\]
We choose one $\rho$ so that $\|zV_\tau R_{0,\tau}(s)\|\leq1/2$
on all these levels, uniformly in $\tau$ and $s\in\gamma_{\rm cut}$.
We write $R_{z,\tau}(s)=(L_{z,\tau}+s\Id)^{-1}$ for the twisted
resolvent. The norm bound gives the convergent series
\[
 (\Id+zV_\tau R_{0,\tau}(s))^{-1}
 =\sum_{n=0}^{\infty}(-zV_\tau R_{0,\tau}(s))^n.
\]
Thus
\begin{equation}\label{Neumann}
 R_{z,\tau}(s)=R_{0,\tau}(s)
                  (\Id+zV_\tau R_{0,\tau}(s))^{-1}
\end{equation}
and is holomorphic on this fixed disk. For two nearby real metric
parameters $\tau$ and $\tau'$, the identity
\[
 R_{z,\tau}-R_{z,\tau'}
 =R_{z,\tau}(L_{z,\tau'}-L_{z,\tau})R_{z,\tau'}
\]
gives continuity $\cH_{j+1}\to\cH_j$ for $0\leq j\leq3$,
locally uniformly in $z$.
For differentiability, \cite[Appendix B.5, (B.6)]{CD} allows a loss
of three levels. The multiplier $V_\tau$ is smooth in $\tau$ on every
chosen level. Applying the same resolvent identity to difference
quotients gives
\[
 \frac{R_{z,\tau'}-R_{z,\tau}}{\tau'-\tau}
 =-R_{z,\tau'}
     \frac{L_{z,\tau'}-L_{z,\tau}}{\tau'-\tau}R_{z,\tau}.
\]
The middle factor converges with a loss of one level. The continuity
just proved controls both outer factors on adjacent levels. Passing to
the limit with a total loss of three levels gives
\begin{equation}\label{metricresolvent}
 \partial_\tau R_{z,\tau}(s)
 =-R_{z,\tau}(s)(\partial_\tau L_{0,\tau}
                    +z\partial_\tau V_\tau)R_{z,\tau}(s)
 \quad\text{as maps }\cH_3\to\cH_0.
\end{equation}
The derivative is jointly continuous and locally uniformly holomorphic
in $z$. All resolvents in this argument stay on the fixed contour;
no individual resonance needs to be differentiated.
Consequently the full-cut projection $\Pi_{z,\tau}$ and its first
moment $\mathsf T_{z,\tau}$, defined by
\begin{equation}\label{xp:projector-moment}
 \begin{aligned}
 \Pi_{z,\tau}&=\frac1{2\pi i}
     \int_{\gamma_{\rm cut}} R_{z,\tau}(s)\,ds \text{ and }
 \mathsf T_{z,\tau}=\frac1{2\pi i}
     \int_{\gamma_{\rm cut}} sR_{z,\tau}(s)\,ds
     =B_{z,\tau}\Pi_{z,\tau},
 \end{aligned}
\end{equation}
have the same regularity. Untwisted ranks persist by
\cite[Appendix B.6]{CD}. To check rank persistence under twisting,
we write $\|\cdot\|_j$ for the operator norm on $\cH_j$ and choose
constants $K_R,K_V>0$ such that, on every chosen level,
\[
 \|R_{0,\tau}(s)\|_j\leq K_R,
 \qquad\|V_\tau R_{0,\tau}(s)\|_j\leq K_V.
\]
The series and the contour length $2\pi\lambda_{\rm cut}$ give
\[
 \|\Pi_{z,\tau}-\Pi_{0,\tau}\|_j
 \leq\lambda_{\rm cut}K_R\frac{|z|K_V}{1-|z|K_V}<1,
\]
after shrinking $\rho$ again. Two projections at distance less than
one have the same finite rank.
We write $N$ for the full cut rank in the chosen degree.
We choose fixed smooth forms $f_1,\ldots,f_N$ and linear functionals
$\varphi_1,\ldots,\varphi_N$ given by pairing with fixed smooth dual
forms. We define the square matrices $F$ and $M_B$ by
\begin{equation}\label{xp:frame-matrices}
 F_{ij}=\varphi_i(\Pi f_j),\qquad (M_B)_{ij}=\varphi_i(\mathsf T f_j)
\end{equation}
We arrange that $F$ is invertible at the origin.
The vectors $\Pi f_j$ then form a basis of the cut near the origin.
The matrix of $B$ in this basis is $F^{-1}M_B$, because
\[
 B\Pi f_j=\sum_{l=1}^N\Pi f_l(F^{-1}M_B)_{lj}.
\]
We write $N_h$ for the horizontal rank in that degree. For the horizontal
cut, we repeat this construction with the projection
$(\Id-\alpha_\tau\wedge\iota_{X_\tau})\Pi$ and with $N_h$ sources
and tests.
Shrinking $\rho,\varepsilon$, we make all full and horizontal frames
invertible throughout the parameter set. Fix $0<\rho_*<\rho$ and put
$\mathcal P=\{(z,\tau):|z|\le\rho_*,\ 0\le\tau\le\varepsilon\}$;
their inverse matrices are uniformly bounded on $\mathcal P$.
The horizontal projection commutes with $\Pi$. We set
$\cU^{-1}=\cU^5=0$. The contact splitting gives
\[
 \dim\ran(\Pi|_{\text{degree }k})
 =\dim\cU^{k-1}+\dim\cU^k\qquad(0\leq k\leq5).
\]
The full ranks therefore determine the horizontal ranks successively,
starting in degree zero. They stay constant. This proves the claimed
matrix regularity.

To prove the symmetry, we work in each horizontal degree separately.
We denote an orbit's transverse symplectic return map by $P_\gamma$,
its length by $\ell_\gamma=\ell_{g_\tau}(\gamma)$, and the integral
of the harmonic form along its projection by
$I_\gamma=\int_{\pi\gamma}\omega$. We use $\Lambda^k$ for the $k$th
exterior power. The degree-$k$ dynamical determinant in a right half-plane is
\[
 \mathscr D_k(s,z,\tau)=\exp\left(-\sum_{\gamma\in\mathcal G^\#_{g_\tau}}
 \sum_{j\geq1}\frac{e^{-sj\ell_\gamma-izjI_\gamma}}{j}
 \frac{\tr(\Lambda^kP_\gamma^j)}{|\det(\Id-P_\gamma^j)|}\right).
\]
We write $\bar\gamma$ for the reversed orbit and use superscript $T$
for matrix transpose. The reversed orbit has
$I_{\bar\gamma}=-I_\gamma$ and return map
similar to $P_\gamma^{-1}$. A symplectic map preserves the transverse
two-form; in particular its determinant is one and its inverse is
similar to its transpose. Thus
\[
 \tr(\Lambda^kP_\gamma^{-j})=\tr(\Lambda^kP_\gamma^j),\qquad
 |\det(\Id-P_\gamma^{-j})|=|\det(\Id-P_\gamma^j)|.
\]
Pairing the two orientations gives
\[
 e^{-izjI_\gamma}+e^{izjI_\gamma}=2\cos(zjI_\gamma).
\]
Hence the determinant satisfies
\[
 \mathscr D_k(s,-z,\tau)=\mathscr D_k(s,z,\tau),\qquad
 \overline{\mathscr D_k(s,z,\tau)}
   =\mathscr D_k(\bar s,\bar z,\tau).
\]
These identities continue meromorphically. Apply
\cite[Section~5, Theorem~4]{DG} to the complex bundle
$\Lambda^k\operatorname{Ann}(X_\tau)\otimes L_{\chi_z}$ with lift $L_{z,\tau}$;
see also \cite[Remark~3]{DGRS}. Here $L_{\chi_z}$ is the flat character line.
The logarithmic derivative has only simple poles, with residues equal to
the algebraic resonance multiplicities; integration gives those
multiplicities as the determinant zero orders, also for complex $z$
and perturbed metrics.
The zeros inside the contour therefore determine the monic cut
polynomial, including multiplicities. This proves the assertion
degree by degree.
\end{proof}

We define operators on horizontal forms by
\[
 D=d_z-\alpha_\tau\wedge L,\qquad W=d\alpha_\tau\wedge.
\]
We write $[B,D]=BD-DB$ for the commutator of $B$ and $D$.
They satisfy
\begin{equation}\label{Cartan}
 [B,D]=0,\qquad D^2=WB.
\end{equation}
Here $D$ is the horizontal part of $d_z$. It does not in general
square to zero. Indeed, for a horizontal form $u$,
\[
 d_zu=Du+\alpha_\tau\wedge Lu,\qquad
 0=d_z^2u=D^2u+WLu,
\]
where the terms containing $\alpha_\tau$ cancel because $D$ and $L$
commute. Thus $D^2u=0$ on zero eigenvectors of $B$.
The two-form $d\alpha_\tau$ is nondegenerate on the four-dimensional
contact plane $\ker\alpha_\tau$. Its wedge maps, called the contact
Lefschetz maps, commute with $B$ and give isomorphisms
\[
 W:\cU^1\longrightarrow\cU^3,\qquad
 W^2:\cU^0\longrightarrow\cU^4.
\]
For small real $\theta\ne0$, pullback acyclicity gives
$H^0(M;\pi^*\chi_\theta)=0$. Hence \cite[Lemma~7.4]{DGRS}
excludes a scalar resonance at zero, so $c\ne0$. For a nonzero scalar
eigenvector $f$ with $Bf=cf$, we have
\[
 D^2f=cWf\ne0,\qquad Df\ne0,\qquad B(Df)=cDf.
\]
Hence $c$ is also a degree-one root. Since the degree-one space has
dimension two and trace $c+t$, its characteristic polynomial is
\begin{equation}\label{P1}
 P_1(s)=(s-c)(s-t),
\end{equation}
first for real nonzero characters and then by holomorphy everywhere.
Using evenness, \eqref{taxis}, and Cauchy's formula, we define quotient
functions $C_c$ and $C_t$ by
\begin{equation}\label{ctfactors}
 c=z^2C_c,\qquad t=t(0,\tau)+z^2C_t,
 \qquad C_c(0,0)=C_t(0,0)=-\kappa/2.
\end{equation}
Here and below the quotients are holomorphic in $z$ and $C^1$ in
$\tau$, locally uniformly in $z$. For any even such function $a_0$ with
$a_0(0,\tau)=0$, and fixed $\rho_*<a<\rho$,
\begin{equation}\label{Cauchyquotient}
 \frac{a_0(z,\tau)}{z^2}
 =\frac1{2\pi i}\int_{|w|=a}
       \frac{a_0(w,\tau)}{w^2(w-z)}\,dw,\qquad |z|<a.
\end{equation}
The removable value and the metric derivative have the same regularity,
by differentiation under the integral. Shrink $\rho_*,\varepsilon$ so that $C_c$ is
nonzero on $\mathcal P$.

When $A_1$ is invertible, $D:\cU^1\to\cU^2$ is injective. Indeed,
\[
 Du=0\ \Longrightarrow\ D^2u=WBu=0
 \ \Longrightarrow\ Bu=0\ \Longrightarrow\ u=0.
\]
Here the second implication uses the isomorphism $W$, and the last
uses invertibility of $A_1$. Since $BD=DB$, the restriction of $B$ to
$D(\cU^1)$ has characteristic polynomial $P_1$. It follows that $P_1$
divides $P_2$.
The remainder in monic polynomial division is holomorphic in $z$
and $C^1$ in $\tau$. For fixed small $\tau>0$, $t(0,\tau)\ne0$,
and \eqref{ctfactors} shows that $c(\cdot,\tau)$ is not identically zero.
Thus $A_1$ is invertible on a dense set of $z$; at $\tau=0$, we use
\eqref{hyper}. The remainder vanishes. We therefore define the monic quadratic
$Q$ and its coefficient functions $e_1,e_2$ by
\begin{equation}\label{P2}
 P_2=(s-c)(s-t)Q(s),\qquad Q(s)=s^2-e_1s+e_2.
\end{equation}
Thus $Q$ contains the two remaining middle-degree roots.
Shrink $\varepsilon$ below the untwisted splitting bound $\varepsilon_0$
established in Appendix~\ref{app:inputs}.
The hyperbolic identities and the untwisted zero multiplicities
in degrees one and two from~\cite[Theorem~3 and Corollary~4.1]{CDDP}
give, respectively,
\begin{equation}\label{qhyper}
 e_1(z,0)=0,\quad e_2(z,0)=-c(z,0)^2,
 \qquad e_1(0,\tau)=e_2(0,\tau)=0.
\end{equation}
For the hyperbolic axis, division of \eqref{hyper} by \eqref{P1} gives
\[
 Q(s,z,0)=(s-c(z,0))(s+c(z,0))=s^2-c(z,0)^2.
\]
For $z=0$ and small $\tau>0$, the horizontal zero ranks are
$(1,1,3,1,1)$. The map $D$ carries the remaining degree-one root
$t(0,\tau)\ne0$ into degree two. The ranks therefore give
\[
 P_1(s,0,\tau)=s(s-t(0,\tau)),\quad
 P_2(s,0,\tau)=s^3(s-t(0,\tau)),\quad Q(s,0,\tau)=s^2.
\]
Continuity gives the same coefficient identities at $\tau=0$.
We apply \eqref{Cauchyquotient} to $e_1$ and call the resulting scalar
quotient $A_e$. Then
\begin{equation}\label{factors}
 e_1=z^2A_e,\qquad A_e(z,0)=0.
\end{equation}
Holomorphy strengthens the two-axis estimate of~\cite{Lawton} to a factor
$z^2\tau$. Define the continuous function, bounded on $\mathcal P$,
\begin{equation}\label{twoaxisfactor}
 E(z,\tau)=\int_0^1\partial_\tau A_e(z,v\tau)\,dv,
 \qquad e_1(z,\tau)=z^2\tau E(z,\tau).
\end{equation}
Since $A_e(z,0)=0$, the fundamental theorem of calculus gives
$A_e(z,\tau)=\tau E(z,\tau)$.

For real $\theta$ and $\tau\geq0$, equations \eqref{taxis} and
\eqref{ctfactors} give
\begin{equation}\label{crossingbalance}
 t(\theta,\tau)=r\tau-\frac\kappa2\theta^2
                   +o(\tau+\theta^2)
\end{equation}
as $(\theta,\tau)\to(0,0)$.
By Lemma~\ref{holom}, shrink $\rho_*,\varepsilon$ so that
$\partial_\tau t(\theta,\tau)>r/2$ for real $(\theta,\tau)\in\mathcal P$.
Thus $t(0,\tau)\geq r\tau/2$, and each crossing below is unique.
Choose $K_*>0$ with $|C_t|\leq K_*$ on $\mathcal P$
and restrict $\theta$ so that $4K_*\theta^2/r<\varepsilon$. Hence
\[
 t(\theta,0)=c_0(\theta)<0,
 \qquad t(\theta,4K_*\theta^2/r)\geq K_*\theta^2>0.
\]
Continuity supplies a crossing $\sigma_\theta$ with
\begin{equation}\label{cross}
 t(\theta,\sigma_\theta)=0,\qquad
 0<\sigma_\theta\leq4K_*\theta^2/r,
 \qquad \frac{\sigma_\theta}{\theta^2}\longrightarrow\frac\kappa{2r}.
\end{equation}
To check the last limit, we use the bound $\sigma_\theta=O(\theta^2)$
and write the root equation as
\[
 0=r\frac{\sigma_\theta}{\theta^2}
       -\frac\kappa2+o(1).
\]
Equation \eqref{twoaxisfactor} now gives
$e_1(\theta,\sigma_\theta)=O(\theta^4)$.
The scalar root $c$ stays negative, since $c=\theta^2C_c$ and
$C_c\to-\kappa/2$.

We now fix a crossing and suppress the parameters. We write $V^k$ for
the horizontal generalized zero space and set $m=\dim V^2$.
Since $c\ne0$, the degree-one zero is simple; Lefschetz gives
\[
 (\dim V^0,\ldots,\dim V^4)=(0,1,m,1,0),
 \qquad L|_{V^1}=L|_{V^3}=0.
\]
We take $\alpha=\alpha_{\sigma_\theta}$ and set $V^{-1}=V^5=0$.
The full generalized zero spaces, with their differential, form the complex
\[
 \cV^k=V^k\oplus\alpha\wedge V^{k-1}\qquad(0\leq k\leq5).
\]
For $u\in V^k$ and $v\in V^{k-1}$, its differential is
\begin{equation}\label{xp:full-zero-complex}
 d_\theta(u+\alpha\wedge v)
   =Du+Wv+\alpha\wedge(Lu-Dv).
\end{equation}
Pullback acyclicity and the resonant de Rham theorem
\cite[Theorem 2.1]{DR} make this full complex exact. We write $D_k$
for $D|_{V^k}$. Since $V^0=V^4=0$, its end terms are
$\cV^1=V^1$ and $\cV^4=\alpha\wedge V^3$. For
$u\in V^3$ and $v\in V^2$, \eqref{xp:full-zero-complex} reduces to
\[
 d_\theta(u+\alpha\wedge v)=-\alpha\wedge D_2v.
\]
In degree one, $L|_{V^1}=0$ gives
\[
 d_\theta|_{\cV^1}=D_1,\qquad
 \ker D_1=\operatorname{im}(d_\theta|_{\cV^0})=0.
\]
In degree four, $\cV^5=0$ gives
\[
 \alpha\wedge D_2(V^2)=\cV^4=\alpha\wedge V^3,
 \qquad D_2(V^2)=V^3.
\]
Thus $D_1$ is injective and $D_2$ is surjective. Also
\[
 D_2D_1=WB|_{V^1}=0.
\]
Consequently
\[
 1=\dim D_1(V^1)\leq\dim\ker D_2=m-1.
\]
The nonzero root $c$ occupies at least one dimension of the rank-four
middle cut. Therefore
\[
 2\leq m\leq3.
\]
We write $\operatorname{ord}_{s=0}$ for the multiplicity of a zero at
$s=0$. Since $c\ne0$ and $t=0$, the polynomial identity gives
\[
 m=\operatorname{ord}_{s=0}P_2
   =1+\operatorname{ord}_{s=0}Q.
\]
It follows that $Q(0)=e_2=0$, so $Q(s)=s(s-e_1)$. Hence
\begin{equation}\label{two}
 e_2=0,\qquad m=3\Longleftrightarrow e_1=0,
 \qquad m=2\Longleftrightarrow e_1\ne0.
\end{equation}

\section{Fixed-cut torsion and the ordinary value}\label{sec:normalization}
We write $Z(s,z,\tau)=Z_{g_\tau,\chi_z}(s)$ and denote by
$G(s,z,\tau)$ the complementary factor in
\cite[(6.4)--(6.5)]{CD}, with equation numbers following author
version~v2. It is holomorphic and nonvanishing in the spectral disk
for each parameter. The full cut in degree $k$ is
$\cU^k\oplus\alpha_\tau\wedge\cU^{k-1}$, with determinant $P_kP_{k-1}$
and we set $P_{-1}=P_5=1$. The contact Lefschetz isomorphisms give
\[
 P_3=P_1,\qquad P_4=P_0.
\]
We define the weighted determinant by
\begin{equation}\label{fulltelescoping}
 \mathfrak D_{\rm cut}(s)
 =\prod_{k=0}^{5}(P_kP_{k-1})^{(-1)^k k}
 =\frac{P_1^2}{P_0^2P_2}
 =\frac{s-t}{(s-c)Q(s)}.
\end{equation}
The two full degrees containing $P_j$ give total exponent
$(-1)^j j+(-1)^{j+1}(j+1)=(-1)^{j+1}$.
In flow dimension five, \cite[(6.5)]{CD} gives
\[
 G=Z\mathfrak D_{\rm cut}.
\]
Therefore
\begin{equation}\label{divisor}
 Z(s,z,\tau)=G(s,z,\tau)
       \frac{(s-c)(s^2-e_1s+e_2)}{s-t}.
\end{equation}
At the crossing, exactness gives $Q(s)=s(s-e_1)$, and
\[
 \frac{Q(s)}{s-t}=\frac{s(s-e_1)}{s}=s-e_1
 \qquad\text{at the crossing}.
\]
This handles both zero multiplicities without root branches~\cite{Lawton}.
Fixed-cut torsion variation controls $G$ along $[0,\sigma_\theta]$.

A dot denotes differentiation with respect to $\tau$. For a
degree-preserving map $T$ of a finite graded space $C^\bullet$, we define
the alternating trace, or supertrace, by
\[
 \str_{C^\bullet}T=\sum_k(-1)^k\tr_{C^k}T.
\]
\begin{lemma}[Variation on an arbitrary fixed cut]\label{cutvariation}
Fix small real $\theta\ne0$. Let $\Pi_{\theta,\tau}$ be any full
resonance cut whose fixed contour remains free of resonances throughout
the parameter interval, and set $C^\bullet_{\theta,\tau}=\ran\Pi_{\theta,\tau}$.
We write $\mathcal T_{\rm fin}$ for its contact-chirality torsion
\cite[Section~3.3]{CD} and define
\[
 H(\theta,\tau)=\str_{C^\bullet_{\theta,\tau}}
       (\Pi_{\theta,\tau}\alpha_\tau\wedge\iota_{\dot X_\tau}).
\]
Then
\[
 \partial_\tau\log\mathcal T_{\rm fin}(\tau)=-H(\theta,\tau).
\]
Nonzero resonances inside the cut are allowed.
For the cut of Section~\ref{sec:spectral}, we also have
\begin{equation}\label{Gvariation}
 \frac{G(0,\theta,\tau)}{G(0,\theta,0)}
 =\exp\left(\int_0^\tau H(\theta,v)\,dv\right).
\end{equation}
For this cut, the finite trace $H$ extends continuously and boundedly
to $\theta=0$. In particular the ratio tends to one uniformly as
$\tau\to0$ and $\theta\to0$ through nonzero real characters.
\end{lemma}
\begin{proof}
We fix $\theta$ and write
$C_\tau=C^\bullet_{\theta,\tau}$ and
$U^k_\tau=C^k_\tau\cap\ker\iota_{X_\tau}$.
We suppress metric subscripts and use Chaubet--Dang's convention
$v\wedge\alpha=(-1)^{k-1}\alpha\wedge v$ for horizontal $(k-1)$-forms.
For a horizontal form $u$, the contact splitting and Lefschetz map are
\[
 C^k_\tau=U^{k-1}_\tau\wedge\alpha_\tau\oplus U^k_\tau,
 \qquad \ell_\tau u=u\wedge d\alpha_\tau,
 \qquad U^{-1}_\tau=U^5_\tau=0.
\]
For $k=0,1,2$, $f\in U^{k-1}$, and $g\in U^k$, we define the contact
chirality $\Gamma$, a degree-reversing involution, by
\begin{equation}\label{xp:contact-chirality}
 \Gamma(f\wedge\alpha+g)
 =\ell^{2-k}g\wedge\alpha+\ell^{3-k}f.
\end{equation}
In higher degrees it is the inverse. Thus
\[
 \Gamma(C^k_\tau)=C^{5-k}_\tau,\qquad\Gamma^2=\Id.
\]
The powers of $\ell$ denote the corresponding Lefschetz maps;
their inverses below act on their ranges.

The generalized zero-resonance part, if present, has the same cohomology as the
twisted de Rham complex of $M$, which is acyclic, by
\cite[Theorem~2.1]{DR}. On the sum of the
nonzero generalized resonance spaces, $L$ is invertible and
\[
 d_\theta(\iota_XL^{-1})+(\iota_XL^{-1})d_\theta=\Id.
\]
Thus the whole cut $C_\tau$ is acyclic. The cut commutes with $d_\theta$,
$\iota_X$, wedging by $\alpha$, and the Lefschetz maps, so
$\Gamma$ preserves it.

Write $\det E$ for the top exterior power of a finite-dimensional space.
The determinant line of the complex is
\[
 \det C_\tau=\bigotimes_{k=0}^5(\det C^k_\tau)^{(-1)^k},
\]
where exponent $-1$ denotes the dual line. Choose volume elements in
degrees $0,1,2$ and their $\Gamma$-partners in degrees $5,4,3$.
Their alternating tensor product, with the sign of~\cite[Section~3.3]{CD},
is choice-independent: paired rescalings cancel. The canonical torsion
isomorphism sends it to $\det H^\bullet(C_\tau)=\mathbb C$ by acyclicity;
its nonzero scalar is $\mathcal T_{\rm fin}(\tau)$.
The smooth-family projector regularity in
\cite[Proposition~7.2 and Appendix~B]{CD} gives differentiable
finite matrices in $\tau$, allowing a loss of anisotropic levels.
We take a reference time $\tau_*$ and write
\[
 \Pi_*=\Pi_{\theta,\tau_*},\qquad
 \dot\Pi_*=(\partial_\tau\Pi_{\theta,\tau})|_{\tau=\tau_*}.
\]
Projection gives the following maps on the fixed space $C=C_{\tau_*}$:
\[
 C=C_{\tau_*},\qquad S_\tau=\Pi_{\theta,\tau}|_C,
 \qquad \mathsf F_\tau=\Pi_*S_\tau:C\to C.
\]
They are invertible near $\tau_*$. Since the projections commute
with $d_\theta$, the differential in the fixed coordinates is unchanged:
\[
 S_\tau^{-1}d_\theta S_\tau=d_\theta|_C.
\]
We define the chirality in these coordinates by
\[
 \widetilde\Gamma_\tau
 =S_\tau^{-1}\Gamma_\tau S_\tau
 =\mathsf F_\tau^{-1}\Pi_*\Gamma_\tau S_\tau.
\]
At $\tau_*$, differentiation of $\Pi_{\theta,\tau}^2=\Pi_{\theta,\tau}$ gives
\[
 \Pi_*\dot\Pi_*\Pi_*=0,
 \qquad \mathsf F_{\tau_*}=\Id,\qquad\dot{\mathsf F}_{\tau_*}=0.
\]
Using $\Pi_*\Gamma_{\tau_*}=\Gamma_{\tau_*}\Pi_*$, we obtain
\[
 \dot{\widetilde\Gamma}_{\tau_*}
 =\Pi_*\dot\Gamma_{\tau_*}|_C
   +\Pi_*\Gamma_{\tau_*}\dot\Pi_*|_C
 =\Pi_*\dot\Gamma_{\tau_*}|_C.
\]
The fixed differential makes the torsion isomorphism constant.
Fix volumes $\upsilon_k\in\det C^k$, $k=0,1,2$.
The varying determinant factors are
$(\widetilde\Gamma_\tau\upsilon_k)^{(-1)^{5-k}}$; their derivatives give
the first equality below.
Differentiating $\widetilde\Gamma^2=\Id$ and pairing complementary degrees
gives the half-supertrace identity~\cite[Proposition~4.9]{BK}, also used in
\cite[Lemma~7.3]{CD}:
\begin{equation}\label{xp:chirality-derivative}
 \begin{aligned}
 \partial_\tau\log\mathcal T_{\rm fin}(\tau)
 &=\sum_{k=0}^{2}(-1)^{k+1}
       \tr_{C^k}(\widetilde\Gamma\dot{\widetilde\Gamma})\\
 &=\tfrac12\str_C(\dot{\widetilde\Gamma}\widetilde\Gamma)
   =\sum_{k=0}^{2}(-1)^{k+1}\tr_{C^k}(\Pi\Gamma\dot\Gamma).
 \end{aligned}
\end{equation}

We define a real function $h$ on $M$, a horizontal two-form $\delta$, and maps
$\mathsf A_0,\mathsf A_1$ on horizontal degrees zero and one by
\begin{equation}\label{xp:chirality-bundle-maps}
 \begin{aligned}
 h&=\alpha(\dot X)=-\dot\alpha(X),
 &\delta&=d\dot\alpha+(\iota_Xd\dot\alpha)\wedge\alpha,\\
 \mathsf A_0u&=2(\ell^2)^{-1}(u\wedge\delta\wedge d\alpha),
 &\mathsf A_1u&=\ell^{-1}(u\wedge\delta).
 \end{aligned}
\end{equation}
In products of operators, $h$ means multiplication by this function.
In particular $\Pi h$ on a cut space means multiplication followed
by projection back to the cut.
To differentiate the chirality, we use
\[
 \partial_\tau(d\alpha_\tau)^j
 =j\,d\dot\alpha_\tau\wedge(d\alpha_\tau)^{j-1},
 \qquad\dot\alpha(X)=-h.
\]
For example, on degree zero, a fixed scalar $u$ gives
\[
 \begin{aligned}
 \dot\Gamma u
 &=2u\,d\dot\alpha\wedge d\alpha\wedge\alpha
      +u(d\alpha)^2\wedge\dot\alpha,\\
 \Gamma\dot\Gamma u&=\mathsf A_0u-hu.
 \end{aligned}
\]
The same differentiation in degrees one and two gives the diagonal
blocks of $\Gamma\dot\Gamma$:
\[
 \begin{array}{c|cc}
 k&U^{k-1}\wedge\alpha&U^k\\ \hline
 0&0&\mathsf A_0-h\\
 1&\mathsf A_0+h&\mathsf A_1-h\\
 2&\mathsf A_1+h&-h
 \end{array}
\]
The off-diagonal blocks have zero compressed trace.
Each $\mathsf A_j$ appears in two adjacent full degrees, so its
terms cancel in the alternating trace. Only the terms involving
the function $h$ remain. To write this cancellation, we set
\[
 a_j=\tr_{U^j}(\Pi\mathsf A_j)\quad(j=0,1),\qquad
 h_j=\tr_{U^j}(\Pi h)\quad(0\leq j\leq4),\qquad
 \gamma_k=\tr_{C^k}(\Pi\Gamma\dot\Gamma).
\]
Then
\begin{equation}\label{xp:chirality-trace-cancellation}
 \begin{aligned}
 \gamma_0&=a_0-h_0,\quad
 \gamma_1=a_0+h_0+a_1-h_1,\quad
 \gamma_2=a_1+h_1-h_2,\\
 -\gamma_0+\gamma_1-\gamma_2&=2h_0-2h_1+h_2.
 \end{aligned}
\end{equation}
For horizontal forms $v$ and $u$, contraction gives
\[
 (\alpha\wedge\iota_{\dot X})(v\wedge\alpha)=h\,v\wedge\alpha,
 \qquad (\alpha\wedge\iota_{\dot X})u
       =\alpha\wedge\iota_{\dot X}u.
\]
The second expression has no horizontal component. Thus the diagonal
of this operator is $h$ on $U^{k-1}\wedge\alpha$ and zero on $U^k$.
Hence
\begin{equation}\label{xp:full-cut-supertrace}
 H=-h_0+h_1-h_2+h_3-h_4=-2h_0+2h_1-h_2.
\end{equation}
The last equality follows from the cut Lefschetz isomorphisms and the
fact that scalar multiplication commutes with them:
\[
 h_3=h_1,\qquad h_4=h_0.
\]
Therefore
\begin{equation}\label{torsionvariation}
 \partial_\tau\log\mathcal T_{\rm fin}(\tau)=-H.
\end{equation}
No restriction on the interior resonances or cut ranks was used.
We now specialize to the cut of Section~\ref{sec:spectral}.
For flow dimension five and the full ranks in \eqref{ranks}, the
sign exponent in \cite[(6.6)]{CD} is
\[
 n_{\rm cut}=\sum_{k=0}^2(-1)^k(3-k)\dim C^k
            =3-6+6=3.
\]
We write $\mathcal T_{\rm dyn}$ for Chaubet--Dang's dynamical torsion.
The finite-cut formula is thus
\begin{equation}\label{dynamicalbridge}
 \mathcal T_{\rm dyn}(\theta,\tau)
 =(-1)^3G(0,\theta,\tau)\mathcal T_{\rm fin}(\tau).
\end{equation}
At a crossing, the zero-complex sign exponent is $m-1$ and
$\operatorname{ord}_{s=0}Z=m-2$. With $\mathcal T_0$ the refined torsion
of the full generalized zero complex, \cite[(6.7)]{CD} gives
$\mathcal T_{\rm dyn}=-Z(0)\mathcal T_0$ for $m=2$, and
$\mathcal T_{\rm dyn}=Z'(0)\mathcal T_0$ for $m=3$, where $Z(0)=0$.
Acyclicity alone does not remove this finite factor; formula~(6.8)
of~\cite{CD} applies when zero is absent from the resonance spectrum.
Deformation invariance~\cite[Theorem~9]{CD} at the fixed acyclic character and
\eqref{torsionvariation} give
\[
 0=\partial_\tau\log\mathcal T_{\rm dyn}
   =\partial_\tau\log G+\partial_\tau\log\mathcal T_{\rm fin}
   =\partial_\tau\log G-H.
\]
Integration proves \eqref{Gvariation}, without choosing a logarithm.

At the trivial character, acyclicity fails but the full cut keeps
its rank. We extend the finite trace $H$ through $\theta=0$;
we do not need to extend either scalar torsion. We set
$\mathcal K_\tau=\alpha_\tau\wedge\iota_{\dot X_\tau}$.
We use the fixed forms $f_j$ and functionals $\varphi_i$ of
\eqref{xp:frame-matrices} to define, in each full degree,
\[
 F_{ij}(z,\tau)=\varphi_i(\Pi_{z,\tau}f_j),\qquad
 K^{\rm var}_{ij}(z,\tau)=\varphi_i
       (\Pi_{z,\tau}\mathcal K_\tau\Pi_{z,\tau}f_j).
\]
The matrices $F$ are uniformly invertible on $\mathcal P$.
The continuity of $K^{\rm var}$
uses precisely the available loss of levels:
\[
 \mathcal H_2\xrightarrow{\Pi}\mathcal H_1
 \xrightarrow{\mathcal K_\tau}\mathcal H_1
 \xrightarrow{\Pi}\mathcal H_0
 \xrightarrow{\varphi_i}\mathbb C.
\]
Writing $F_k,K^{\rm var}_k$ for the matrices in degree $k$, we obtain the
continuous extension
\[
 H(z,\tau)=\sum_{k=0}^5(-1)^k
          \tr(F_k(z,\tau)^{-1}K^{\rm var}_k(z,\tau)).
\]
Choose $K_H\geq\sup_{\mathcal P}|H|$. Equation \eqref{Gvariation}
then yields
\[
 \left|\frac{G(0,\theta,\tau)}{G(0,\theta,0)}-1\right|
 \leq e^{K_H\tau}-1.
\]
\end{proof}

\begin{proof}[Proof of Theorem~\ref{mainresult}]
The four inputs established above are
\[
 \begin{aligned}
 c_0(\theta)&=-\tfrac\kappa2\theta^2+O(\theta^4),
 &\partial_\tau t(0,0)&=r>0,\\
 e_1(z,\tau)&=z^2\tau E(z,\tau),
 &\frac{G(0,\theta,\sigma_\theta)}{G(0,\theta,0)}&=1+O(\theta^2).
 \end{aligned}
\]
At the hyperbolic endpoint, \eqref{hyper}, \eqref{qhyper}, and
\eqref{divisor} give
\begin{equation}\label{hypervalue}
 Z_{\rm hyp}(\theta):=Z(0,\theta,0)
       =-G(0,\theta,0)c_0(\theta)^2\ne0.
\end{equation}
We fix a crossing $(\theta,\sigma_\theta)$ from \eqref{cross}.
Equation \eqref{two} gives $Q(s)=s(s-e_1)$. We first specialize the
parameters and then cancel the common $s$ in
\eqref{divisor}. This gives
\begin{equation}\label{specialized}
 \begin{aligned}
 Z(s,\theta,\sigma_\theta)
 &=G(s,\theta,\sigma_\theta)(s-c)(s-e_1),\\
 Z(0,\theta,\sigma_\theta)&=G(0,\theta,\sigma_\theta)c e_1.
 \end{aligned}
\end{equation}
Here $c,e_1$ are evaluated at the crossing. If $m=3$, this has
a simple zero at zero; if $m=2$, its value there is nonzero.
In either case it is holomorphic and
\begin{equation}\label{exactratio}
 \frac{Z(0,\theta,\sigma_\theta)}{Z_{\rm hyp}(\theta)}
 =-\frac{G(0,\theta,\sigma_\theta)}{G(0,\theta,0)}
   \frac{c(\theta,\sigma_\theta)}{c_0(\theta)}
   \frac{e_1(\theta,\sigma_\theta)}{c_0(\theta)}.
\end{equation}
The first factor is $1+O(\theta^2)$ by Lemma~\ref{cutvariation} and
$\sigma_\theta=O(\theta^2)$. The bounded metric derivative gives
\[
 \frac{C_c(\theta,\sigma_\theta)}{C_c(\theta,0)}-1
 =\frac{\sigma_\theta}{C_c(\theta,0)}
   \int_0^1\partial_\tau C_c(\theta,v\sigma_\theta)\,dv
 =O(\theta^2).
\]
Together with \eqref{twoaxisfactor}, this proves
\begin{equation}\label{coefficientratios}
 \begin{aligned}
 \frac{c(\theta,\sigma_\theta)}{c_0(\theta)}
 &=\frac{C_c(\theta,\sigma_\theta)}{C_c(\theta,0)}=1+O(\theta^2) \text{ and }
 \frac{e_1(\theta,\sigma_\theta)}{c_0(\theta)}
 =\frac{\sigma_\theta E(\theta,\sigma_\theta)}{C_c(\theta,0)}
   =O(\theta^2).
 \end{aligned}
\end{equation}
To obtain the leading coefficient, we divide the second ratio by $\theta^2$:
\[
 \frac{e_1(\theta,\sigma_\theta)}{\theta^2c_0(\theta)}
 =\frac{\sigma_\theta}{\theta^2}
   \frac{E(\theta,\sigma_\theta)}{C_c(\theta,0)}
 \longrightarrow\frac\kappa{2r}\frac{E(0,0)}{-\kappa/2}
 =-\frac{E(0,0)}r.
\]
The minus sign in \eqref{exactratio} therefore gives
\begin{equation}\label{quantitativeratio}
 \frac{Z(0,\theta,\sigma_\theta)}{Z_{\rm hyp}(\theta)}
 =\frac{E(0,0)}{r}\theta^2+o(\theta^2).
\end{equation}
We do not evaluate $E(0,0)$. If it vanishes, the ratio is
$o(\theta^2)$. An actual zero at a crossing requires
$e_1(\theta,\sigma_\theta)=0$, equivalently $m=3$ by~\eqref{two}.
Fried's hyperbolic theorem \cite{Fried}, in the complex unitary
form stated in \cite[Theorem 3]{DGRS}, gives
\[
 |Z_{\rm hyp}(\theta)|=T_{\chi_\theta}(\Sigma)^2.
\]
In particular, the torsion-normalized value satisfies
\begin{equation}\label{quantitativebound}
 \frac{|Z_{g_{\sigma_\theta},\chi_\theta}(0)|}
      {T_{\chi_\theta}(\Sigma)^2}=O(\theta^2).
\end{equation}
For fixed sufficiently small real $\theta\ne0$ and
$\tau\in[0,\varepsilon]\setminus\{\sigma_\theta\}$, we have $ct\ne0$.
Equation~\eqref{P1} and the contact splitting exclude a full degree-one
zero resonance. Pullback acyclicity and \cite[Proposition~7.3]{DGRS}
then exclude zero in every degree. Thus \cite[Theorem~9 and (6.8)]{CD}
give $Z(0,\theta,\tau)=Z_{\rm hyp}(\theta)$, proving
\eqref{Friedidentity} at every other metric in this small family.
Acyclic unitary analytic torsion is metric-independent on the
odd-dimensional base. Take sufficiently small real $\theta_j$ in the
parameter range fixed above and set
\[
 0\ne\theta_j\longrightarrow0,\qquad
 g_j=g_{\sigma_{\theta_j}},\qquad\chi_j=\chi_{\theta_j}.
\]
This proves the theorem. Each metric is real analytic in the hyperbolic atlas of $g_0$,
since the fixed conformal function $b$ is real analytic.
The metrics converge smoothly because
$\sigma_{\theta_j}\to0$.
\end{proof}

\appendix
\section{Geometric and spectral inputs}\label{app:inputs}

\paragraph{The base and its characters.}
In this appendix, we write $b_k=b_k(\Sigma)$.
The subscript on $\Id$ gives the matrix size.
We define a symmetric matrix $A_{\rm hom}$, an integral homology action $T_{\rm hom}$
on a genus-two surface, and its characteristic polynomial $p$ by
\[
 A_{\rm hom}=\begin{pmatrix}2&1\\1&-1\end{pmatrix},\qquad
 T_{\rm hom}=\begin{pmatrix}A_{\rm hom}&-\Id_2\\\Id_2&0\end{pmatrix}.
\]
The symmetry of $A_{\rm hom}$ makes $T_{\rm hom}$ preserve the standard
symplectic form. Direct calculation gives
\[
 \begin{aligned}
 p(x)&:=\det(x\Id_4-T_{\rm hom})
 =\det\bigl((x^2+1)\Id_2-xA_{\rm hom}\bigr)\\
 &=(x^2-2x+1)(x^2+x+1)-x^2
 =x^4-x^3-x^2-x+1.
 \end{aligned}
\]
We write $\Phi_5$ for the fifth cyclotomic polynomial. Reduction modulo
two gives
\[
 p(x)\bmod2=x^4+x^3+x^2+x+1=\Phi_5(x)\bmod2.
\]
This polynomial is irreducible over the field with two elements because
the order of $2$ modulo $5$ is $4$. Hence $p$ is irreducible over
$\mathbb Q$. Also
\[
 p(1)=-1,\qquad p(2)=3.
\]
Thus $p$ has a real root greater than one. Since $p$ is irreducible,
a root of unity would force $p$ to equal its cyclotomic minimal polynomial,
contradicting that real root. Hence $p$ has no roots of unity as zeros.
Finally, for any polynomial
$q$ and integer $k>1$,
\[
 p'(0)=-1,\qquad
 \left.\frac{d}{dx}q(x^k)\right|_{x=0}=0.
\]
Therefore $p(x)$ cannot equal $q(x^k)$.
The mapping-class action onto the integral symplectic group and
\cite[Proposition~1]{MS} therefore give an orientation-preserving
pseudo-Anosov mapping class inducing $T_{\rm hom}$.
Its mapping torus $\Sigma$ is hyperbolic by \cite{Thurston}.
The homology sequence of the mapping torus, called the Wang sequence, gives
\[
 b_1(\Sigma)=1+\dim\ker(T_{\rm hom}-\Id_4)=1,
\]
because $p(1)=-1$ makes $T_{\rm hom}-\Id_4$ invertible.
We choose a hyperbolic metric $g_0$ and a nonzero real harmonic
one-form $\omega$.
The hyperbolic charts give $\Sigma$ a real analytic atlas.

We use the adjoints and nonnegative twisted Hodge Laplacians defined
in Section~\ref{sec:spectral}.
All norms below are the corresponding $L^2$ norms.
We choose $\eta>0$ below the first positive untwisted scalar and
one-form eigenvalues.
Their spectral projections onto $[0,\eta)$ each retain rank one
for small $\theta$.
A loop with nonzero $\omega$-period proves that $\chi_\theta$ is
nontrivial for all sufficiently small nonzero $\theta$.
The scalar eigenvalue $\mu_0(\theta)$ in this one-dimensional space is then positive:
a harmonic scalar section would be parallel.
For a normalized eigenfunction $f_\theta$,
\[
 \|d_\theta f_\theta\|_2^2=\mu_0(\theta)>0,\qquad
 \Delta_{\mathrm H,\theta}^{(1)}d_\theta f_\theta
   =\mu_0(\theta)d_\theta f_\theta.
\]
Since the one-form spectral space below $\eta$ has dimension one,
$d_\theta f_\theta$ spans it. Its eigenvalue is positive, so
\[
 \ker\Delta_{\mathrm H,\theta}^{(0)}
 =\ker\Delta_{\mathrm H,\theta}^{(1)}=0,
 \qquad H^0(\Sigma;\chi_\theta)=H^1(\Sigma;\chi_\theta)=0.
\]
Unitary Hodge duality gives the remaining degrees:
\[
 H^k(\Sigma;\chi_\theta)=0\qquad(0\leq k\leq3).
\]
The unit tangent bundle is an oriented sphere bundle with fiber the
two-sphere $S^2$. For all cohomological degrees $p,q$, the spectral
sequence for this bundle has terms
\[
 H^p(\Sigma;\chi_\theta)\otimes H^q(S^2;\mathbb C)=0.
\]
It follows that
\[
 H^\bullet(M;\pi^*\chi_\theta)=0.
\]
For small nonzero real $\theta$, the same argument proves a
uniform gap $\eta$ on coclosed one-forms. The only small one-form
eigenspace consists of exact forms, and its eigenvalue is positive.
At $\theta=0$, this gap holds after removing the harmonic
one-form space.

We write $d\operatorname{vol}_{g_0}$ for the Riemannian volume form.
We set $V_\Sigma=\operatorname{vol}_{g_0}(\Sigma)$ and use
$\kappa=\|\omega\|_2^2/V_\Sigma$ as in the main proof.
In the next formula, $\langle\cdot,\cdot\rangle$ is the pointwise
inner product from $g_0$. We define scalar operators $K_1,K_2$ by
$K_1f=-2i\langle\omega,df\rangle$ and $K_2f=|\omega|^2f$
for a smooth function $f$. Harmonicity of $\omega$ gives
\[
 \Delta_{\mathrm H,\theta}^{(0)}
 =\Delta_{\mathrm H,0}^{(0)}+\theta K_1+\theta^2K_2.
\]
We take the normalized constant $f_0=V_\Sigma^{-1/2}$. Then
\[
 K_1f_0=0,\qquad\mu_0'(0)=0.
\]
Differentiating the scalar eigenvalue equation at zero and normalizing
$f_\theta$ so that $f'_0$ is orthogonal to $f_0$ gives
\[
 \Delta_{\mathrm H,0}^{(0)}f'_0=-K_1f_0+\mu_0'(0)f_0=0,
 \qquad f'_0=0.
\]
The second derivative therefore gives
\[
 \frac{\mu_0''(0)}2
 =\langle f_0,K_2f_0\rangle_{L^2}
 =\frac1{V_\Sigma}\int_\Sigma|\omega|^2\,d\operatorname{vol}_{g_0}
 =\kappa.
\]
Simple-eigenvalue perturbation gives an analytic eigenvalue branch,
and complex conjugation gives $\mu_0(-\theta)=\mu_0(\theta)$. Thus
\begin{equation}\label{app:scalarquadratic}
 \mu_0(\theta)=\kappa\theta^2+O(\theta^4),\qquad
 c_0(\theta)=-\frac{\mu_0(\theta)}{1+\sqrt{1-\mu_0(\theta)}}
 =-\frac\kappa2\theta^2+O(\theta^4).
\end{equation}

\paragraph{The conformal sign.}
Here $z=\tau=0$, so $X=X_0$, $\alpha=\alpha_0$, and $\Delta_{g_0}$ is the
nonpositive scalar Laplacian used in \cite{CDDP}.
In the Levi--Civita splitting of $TM$, we write tangent vectors as
$\xi=(\xi_H,\xi_V)$ and $\eta=(\eta_H,\eta_V)$, with horizontal
and vertical components. The cross product uses the metric and
orientation on $\Sigma$. We define the tangent rotation and two-form
\[
 \mathsf I_{(x,v)}(\xi_H,\xi_V)
   =(v\times\xi_V,v\times\xi_H),\qquad
 \psi(\xi,\eta)=d\alpha(\mathsf I\xi,\eta).
\]
We write $J(x,v)=(x,-v)$ for time reversal and use $\pi_*$ for
fiber integration with
the convention of \cite{CDDP}.
Let $\cC$ be the closed horizontal zero one-form space and
$\cC_\psi=\{v\circ\mathsf I:v\in\cC\}$ its complementary summand
from~\cite[Lemma~3.9]{CDDP}. The pushforward isomorphism in
\cite[Lemma~3.17]{CDDP} commutes with conjugation; its inverse image
$u\in\cC_\psi$ of the real form $\omega$ is therefore real and satisfies
$\pi_*(d\alpha\wedge u)=\omega$.
Then $du\ne0$ and $u\circ\mathsf I$ is closed.
Set $u_*=J^*u$ and $N_\omega=\|\omega\|_2^2$.
The results in~\cite[Lemma~3.17 and Corollary~5.1]{CDDP} identify
the pushforward of $u_*$:
\[
 \pi_*(d\alpha\wedge u_*)=\omega.
\]
The forms $u,u_*$ may be distributions; $\operatorname{WF}$ denotes wavefront set.
Their outgoing and incoming conditions give
\[
 \operatorname{WF}(du)\cap\bigl(-\operatorname{WF}(du_*)\bigr)=\varnothing.
\]
Thus their product is defined; see~\cite[Section~2.3]{CDDP}.
We write $\mathcal D'(\Sigma;\mathbb R)$ for real distributions on
$\Sigma$. We define a distribution $F_{\rm dist}$ and a smooth function $f$ by
\[
 \pi_*(\alpha\wedge du\wedge du_*)
       =F_{\rm dist}\,d\operatorname{vol}_{g_0},\qquad
 F_{\rm dist}\in\mathcal D'(\Sigma;\mathbb R),\qquad
 f=\Delta_{g_0}|\omega|^2.
\]
We write $\mathbb H^3$ for hyperbolic three-space and
$d\operatorname{vol}_{\mathbb H^3}$ for its volume form. For a function $h$
on $\Sigma$, we write $\widetilde h$ for its lift to $\mathbb H^3$.
We write $d_{\mathbb H^3}$ for its distance,
and define $Q_4$ by
\begin{equation}\label{app:Q4definition}
 \widetilde{Q_4h}(x)=\int_{\mathbb H^3}
 \frac{\widetilde h(y)}{\cosh^4 d_{\mathbb H^3}(x,y)}
       \,d\operatorname{vol}_{\mathbb H^3}(y).
\end{equation}
The kernel gives a symmetric smooth kernel on the compact quotient.
Thus $Q_4$ extends to distributions by duality and maps them to smooth functions.
By~\cite[Theorem~5 and Proposition~A.1]{CDDP}, we have
\[
 Q_4F_{\rm dist}=-f/6,\qquad f\ne0.
\]
We take $b=Q_4f$ as our conformal direction.
This function is real analytic. Indeed, after a hyperbolic isometry,
any point can be taken to $(1,0)$ in the hyperboloid model.
We use absolute values for scalars and Euclidean norms for real or
complex vectors. The dot below denotes the bilinear Euclidean product.
The model is
\[
 \mathbb H^3=\{(y_0,y')\in\mathbb R\times\mathbb R^3:
                  y_0>0,\ y_0^2-|y'|^2=1\}.
\]
For a complex point $w=(w_0,w')\in\mathbb C\times\mathbb C^3$
near $(1,0)$, the kernel in
\eqref{app:Q4definition} extends to $(w_0y_0-w'\cdot y')^{-4}$.
Since $|y'|\leq y_0$, the condition
$|w_0-1|+|w'|<1/2$ gives
\[
 \begin{aligned}
 |w_0y_0-w'\cdot y'|
 &\geq y_0-|w_0-1|y_0-|w'||y'|\\
 &\geq(1-|w_0-1|-|w'|)y_0>y_0/2.
 \end{aligned}
\]
The lifted function $\widetilde f$ is bounded. In polar coordinates,
with distance variable $\varrho$, we have $y_0=\cosh\varrho$ and
\[
 \int_{\mathbb H^3}y_0^{-4}\,d\operatorname{vol}_{\mathbb H^3}
 =4\pi\int_0^\infty\frac{\sinh^2\varrho}{\cosh^4\varrho}\,d\varrho
 =\frac{4\pi}{3}<\infty.
\]
Thus the integral extends holomorphically near $(1,0)$.
Its restriction to the real hyperboloid is real analytic.
Applying this argument at every point proves the claim on $\Sigma$.
It follows that every metric $g_\tau=e^{-2\tau b}g_0$ is real analytic
in the hyperbolic atlas of $g_0$.

With contact orientation on $M$, the fiber convention is
$\pi_*(\alpha\wedge(d\alpha)^2)=-8\pi\,d\operatorname{vol}_{g_0}$.
Below, $\langle F_{\rm dist},h\rangle$ is the pairing of a distribution
with a smooth test function $h$, using the volume form of $g_0$.
For two smooth functions it is the $L^2(g_0)$ pairing.
We define the scalar $S$ by the following integral.
The identity in~\cite[(5.6)]{CDDP} gives
\[
 \begin{aligned}
 S&:=\int_M\pi^*b\,\alpha\wedge du\wedge du_*=-\langle F_{\rm dist},Q_4f\rangle
   =-\langle Q_4F_{\rm dist},f\rangle
   =\frac16\|f\|_2^2>0.
 \end{aligned}
\]
For outgoing and incoming one-forms $v,v_*$, we define the pairing
$\mathcal B$ by
$\mathcal B(v,v_*)=\int_M\alpha\wedge d\alpha\wedge v\wedge v_*$.
We set $u_{\rm cl}=u\circ\mathsf I$ and use brackets for de Rham
classes. The identities in
\cite[Lemmas~2.10 and 3.5, (3.57), (3.67), and the calculation preceding (3.59)]{CDDP}
give
\[
 [u_{\rm cl}]=-[\pi^*\omega]/(4\pi),\qquad
 \mathcal B(u_{\rm cl},J^*u_{\rm cl})=-N_\omega/(4\pi),\qquad
 \mathcal B(u,u_*)=N_\omega/(4\pi).
\]
The same identities give $\pi_*(\psi\wedge u_{\rm cl})=\omega$.
The closed and nonclosed degree-one blocks each have dimension one,
and their pairing is diagonal by~\cite[(3.55)--(3.57)]{CDDP}.
We write $P'_{\rm cut}$ and $B'_{\rm cut}$ for the metric derivatives
at zero compressed to the degree-one cut; here $P=-iL=iB$.
The identification $F_\tau$ of the unit tangent bundles gives
\[
 \alpha_\tau=e^{-\tau\pi^*b}\alpha_0,
 \qquad\iota_X\partial_\tau\alpha_\tau|_0=-b\circ\pi.
\]
The identities in~\cite[Lemma~4.4, (4.22), and (4.38)]{CDDP} therefore give
\[
 \mathcal B(P'_{\rm cut}u,u_*)=iS.
\]
In the ordered basis $(u_{\rm cl},u)$, the entries in the closed row
and column vanish because $du_{\rm cl}=d(J^*u_{\rm cl})=0$.
The remaining diagonal entry is
\[
 \frac{\mathcal B(P'_{\rm cut}u,u_*)}{\mathcal B(u,u_*)}
 =\frac{iS}{N_\omega/(4\pi)}
 =i\,\frac{4\pi S}{N_\omega}
 =i\,\frac{2\pi}{3N_\omega}\|f\|_2^2.
\]
Consequently
\begin{equation}\label{app:positivevariation}
 P'_{\rm cut}=\operatorname{diag}(0,ir),\qquad
 B'_{\rm cut}=\operatorname{diag}(0,r),\qquad
 r=\frac{2\pi}{3N_\omega}\|f\|_2^2>0.
\end{equation}
Reversing $b$ reverses $S$ and $B'_{\rm cut}$; our choice already gives
$r>0$. Constants span the rank-one scalar cut, so $c(0,\tau)=0$.
The finite-cut trace identity~\cite[Lemma~A.2]{CD} gives
$\partial_\tau t(0,0)=\tr B'_{\rm cut}=r$; since $t(0,0)=0$,
this proves \eqref{taxis}.
For $b_1=1$, this is the nondegeneracy condition
in~\cite[Theorem~3]{CDDP}. Thus
\cite[Theorem~3 and Corollary~4.1]{CDDP} give $\varepsilon_0>0$ such that
the horizontal untwisted zero-resonance ranks are $(1,1,3,1,1)$
for $0<\tau<\varepsilon_0$.

\paragraph{The hyperbolic clusters.}
We verify~\eqref{app:clusterpolynomials} by factoring the horizontal
dynamical determinants near zero.
We write $SO(m)$ for the special orthogonal group in dimension $m$.
For $l=0,1,2$ and integers $p\geq0$, we denote by $\nu_l$ and $\sigma_p$
the $SO(2)$ representations on
$\Lambda^l\mathbb C^2$ and trace-free symmetric tensors of order
$p$, respectively.
Here $\mathbf1$ is the trivial representation; weight $p$ means that
rotation through $\varphi$ acts by $e^{ip\varphi}$.
We have
\[
 \nu_0=\nu_2=\sigma_0=\mathbf1,\qquad
 \sigma_1=\nu_1,\qquad
 \nu_1\otimes\nu_1=\mathbf1^{\oplus2}\oplus\sigma_2.
\]
The weights of $\sigma_1$ are $\pm1$, and those of $\sigma_2$ are $\pm2$.
For a complex variable $a$ and an $SO(2)$ representation $\nu$, we write
$Z^D_\nu(a)$ for the Selberg function in~\cite{DGRS}, with character
$\chi_\theta$. The sign $\doteq$ means equality up to a holomorphic
factor with no zeros near the spectral point in question.
Write $\Det_\zeta$ for the zeta-regularized determinant; graded
determinants are numerator determinants divided by denominator determinants.
With the centered Selberg variable $q=a-1$, the graded determinant
formula in~\cite[Proposition~4.7]{BO} gives
\begin{align}\label{app:determinants}
 Z^D_{\mathbf1}(a)&\doteq
 \Det_\zeta(\Delta_{\mathrm H,\theta}^{(0)}-1+(a-1)^2),\\
 Z^D_{\nu_1}(1+q)&\doteq
 \frac{\Det_\zeta(\Delta_{\mathrm H,\theta}^{(1)}+q^2)}
      {\Det_\zeta(\Delta_{\mathrm H,\theta}^{(0)}+q^2)}.
\end{align}
Each weight pair occurs once.
The exact one-form spectrum cancels the scalar denominator.
At $\theta=0$, the harmonic contribution cancels as well, since
$b_0=b_1=1$.
The coclosed gap above therefore makes the centered weight-one
factor holomorphic and nonvanishing on a fixed small disk.

The weight-two factor uses the formal difference
$\operatorname{Sym}^2_0\mathbb C^3-\Lambda^1\mathbb C^3$ of $SO(3)$
representations, with $\operatorname{Sym}^2_0$ denoting trace-free symmetric
tensors. Below they denote the associated bundles twisted by $\chi_\theta$;
$\nabla_\theta$ is the tensor connection and $\nabla_\theta^*$ its $L^2(g_0)$ adjoint.
Appendix~B of~\cite{DGRS} and the same determinant formula give, at $a=s+2$,
the numerator and denominator operators
\[
 \begin{gathered}
 \nabla_\theta^*\nabla_\theta-2+2s+s^2
       \quad\hbox{on }\operatorname{Sym}^2_0\mathbb C^3 \text{ and }
 \nabla_\theta^*\nabla_\theta+2+2s+s^2
       \quad\hbox{on }\Lambda^1\mathbb C^3.
 \end{gathered}
\]
The Bochner bound used in the proof of~\cite[Proposition~7.7]{DGRS} is
$\nabla_\theta^*\nabla_\theta\ge3$ on $\operatorname{Sym}^2_0$;
a flat unitary twist changes no curvature term.
At $s=0$ both operators are strictly positive. Adding the scalar
$2s+s^2$ keeps them invertible when $|2s+s^2|<1/2$.

\begin{samepage}
In~\cite[(7.6)]{DGRS}, the radial index is $j\ge0$, the exterior index
is $0\le l\le k$, and the symmetric degree is $p\ge0$.
The representation is $\nu_l\otimes\nu_{k-l}\otimes\sigma_p$
and its shift is $\mathfrak h=2(j-l)+p+2+k$.
For $k=0,1,2$, the complete list with $\mathfrak h\le2$ is
\[
\begin{array}{ccccl}
 k&l&p&j&\text{factor}\\\hline
 0&0&0&0&Z^D_{\mathbf1}(s+2)\\
 1&1&0&0&Z^D_{\nu_1}(s+1)\\
 1&1&1&0&Z^D_{\nu_1\otimes\nu_1}(s+2)\\
 2&2&0&0&Z^D_{\mathbf1}(s)\\
 2&2&0&1&Z^D_{\mathbf1}(s+2)\\
 2&2&1&0&Z^D_{\nu_1}(s+1)\\
 2&2&2&0&Z^D_{\sigma_2}(s+2)\\
 2&1&0&0&Z^D_{\nu_1\otimes\nu_1}(s+2).
\end{array}
\]
\end{samepage}
We now bound the infinite tail. We write $\ell_*>0$ for the systole,
$q_*=e^{-\ell_*}$, and $\ell_\gamma=\ell_{g_0}(\gamma)$.
For fixed $k$, let $I_k$ be the triples $(j,l,p)$ with $\mathfrak h\ge3$, and put
$\mu_{klp}=\nu_l\otimes\nu_{k-l}\otimes\sigma_p$.
For $0<w<1$, the representation dimensions give
\[
 \sum_{j,p\ge0}\dim\sigma_p\,w^{2j+p}=(1-w)^{-2},
 \qquad\sum_{l=0}^k\dim\nu_l\dim\nu_{k-l}=\binom4k\le6.
\]
Since $2+k-2l\ge0$, for a real variable $t\ge\ell_*$ we obtain
\[
 \sum_{I_k}\dim\mu_{klp}\,e^{-\mathfrak h t}
 \le e^{-3(t-\ell_*)}\sum_{I_k}\dim\mu_{klp}\,e^{-\mathfrak h\ell_*}
 \le\frac{6e^{3\ell_*}}{(1-q_*)^2}e^{-3t}.
\]
For $n\ge1$, we denote the contracting determinant in
\cite[(7.5)]{DGRS} by $d_{\gamma,n}$. It satisfies
\[
 (1-q_*)^2\le d_{\gamma,n}\le(1+e^{-n\ell_\gamma})^2\le4.
\]
For real $\theta$, the character $\chi_\theta$ is unitary.
Thus its Euler-series logarithms satisfy, for $0<\delta_s<1/2$,
\[
 \sum_{I_k}\sup_{\substack{\operatorname{Re}s\ge-\delta_s\\\theta\in\mathbb R}}
 |\log Z^D_{\mu_{klp}}(s+\mathfrak h)|
 \le\frac{6e^{3\ell_*}}{(1-q_*)^4}
 \sum_{\gamma\in\mathcal G_{g_0}^\#}\sum_{n\ge1}
 \frac{e^{-(3-\delta_s)n\ell_\gamma}}{n}<\infty.
\]
The last sum is bounded by four times the convergent scalar Euler
logarithm at $3-\delta_s>2$~\cite[Section~7.3]{DGRS}.
The exponential of this uniformly convergent logarithm sum is holomorphic
and nonvanishing on $\operatorname{Re}s>-\delta_s$.
This half-plane overlaps $\operatorname{Re}s>2$, where
\cite[(7.6)]{DGRS} holds, so the factorization continues to zero.
For the horizontal dynamical determinants $\mathscr D_k$
of Lemma~\ref{holom}, the table therefore gives
\[
 \begin{aligned}
 \mathscr D_0&\doteq Z^D_{\mathbf1}(s+2),\\
 \mathscr D_1&\doteq
 Z^D_{\nu_1}(s+1)Z^D_{\mathbf1}(s+2)^2Z^D_{\sigma_2}(s+2)
 \doteq Z^D_{\mathbf1}(s+2)^2,\\
 \mathscr D_2&\doteq
 Z^D_{\mathbf1}(s)Z^D_{\mathbf1}(s+2)^3
 Z^D_{\nu_1}(s+1)Z^D_{\sigma_2}(s+2)^2\doteq Z^D_{\mathbf1}(s)Z^D_{\mathbf1}(s+2)^3.
 \end{aligned}
\]
We remove only the weight-one and weight-two factors already shown
to be holomorphic and nonvanishing. The scalar factors give the equations
\[
 \begin{aligned}
 \mu_0+s^2+2s&=(s-c_0)(s+2+c_0)\text{ and }
 \mu_0+s^2-2s=(s+c_0)(s-2-c_0).
 \end{aligned}
\]
Their simple roots near zero are $c_0$ and $-c_0$, respectively;
the other roots tend to $-2$ and $2$.
All other scalar eigenvalues stay separated from zero.
The multiplicity statement following (7.3) in~\cite{DGRS} identifies
the determinant orders with algebraic resonance multiplicities, giving
\begin{equation}\label{app:clusterpolynomials}
 P_0(s)=s-c_0,\qquad P_1(s)=(s-c_0)^2,\qquad
 P_2(s)=(s-c_0)^3(s+c_0).
\end{equation}
At $\theta=0$, \cite[Proposition~7.7]{DGRS} gives the same horizontal
ranks $(1,2,4,2,1)$, with the last two degrees supplied by contact
Lefschetz.



\begin{thebibliography}{DGRS}

\bibitem[BK]{BK} M. Braverman and T. Kappeler,
\emph{Refined analytic torsion as an element of the determinant line},
Geom. Topol. \textbf{11} (2007), 139--213,
\href{https://doi.org/10.2140/gt.2007.11.139}{doi:10.2140/gt.2007.11.139}.

\bibitem[DG]{DG} S. Dyatlov and C. Guillarmou,
\emph{Pollicott--Ruelle resonances for open systems},
Ann. Henri Poincar\'e \textbf{17} (2016), no.~11, 3089--3146,
\href{https://doi.org/10.1007/s00023-016-0491-8}{doi:10.1007/s00023-016-0491-8},
\href{https://math.mit.edu/~dyatlov/files/2014/rnc.pdf}{author version},
Section~5, Theorem~4.

\bibitem[L]{Lawton} B. Lawton,
\emph{Resonance-cluster analysis for a proposed counterexample to Fried's conjecture},
September 2026,
\href{https://github.com/blawton/fried-conj-lean/tree/610d659ea15af869553f6e3e79c432880d9b0ab0}{github}.
See \href{https://github.com/blawton/fried-conj-lean/blob/610d659ea15af869553f6e3e79c432880d9b0ab0/docs/counterexample_intuitive_view.tex}{the companion account}
and \href{https://github.com/blawton/fried-conj-lean/blob/610d659ea15af869553f6e3e79c432880d9b0ab0/B1s/fried_crossing_purezeros_9_08.lean}{the two-variable vanishing lemma}.

\bibitem[RS]{RS} D. B. Ray and I. M. Singer,
\emph{$R$-torsion and the Laplacian on Riemannian manifolds},
Adv. Math. 7 (1971), no. 2, 145--210,
\href{https://doi.org/10.1016/0001-8708(71)90045-4}{doi:10.1016/0001-8708(71)90045-4}.

\bibitem[Z]{Zw} M. Zworski,
\emph{Mathematical study of scattering resonances},
Bull. Math. Sci. 7 (2017), 1--85,
\href{https://doi.org/10.1007/s13373-017-0099-4}{DOI},
Conjecture 9.

\bibitem[HKS]{HKS} C. Hadfield, S. Kandel, and M. Schiavina,
\emph{Ruelle zeta function from field theory},
Ann. Henri Poincar\'e \textbf{21} (2020), no.~12, 3835--3867,
\href{https://doi.org/10.1007/s00023-020-00964-8}{doi:10.1007/s00023-020-00964-8},
Conjecture~21, equation~(35).

\bibitem[F95]{Fried95} D. Fried,
\emph{Meromorphic zeta functions for analytic flows},
Comm. Math. Phys. \textbf{174} (1995), no.~1, 161--190,
\href{https://doi.org/10.1007/BF02099469}{doi:10.1007/BF02099469};
\href{https://projecteuclid.org/journals/communications-in-mathematical-physics/volume-174/issue-1/Meromorphic-zeta-functions-for-analytic-flows/cmp/1104275099.pdf}{primary PDF}.

\bibitem[F86]{Fried} D. Fried,
\emph{Analytic torsion and closed geodesics on hyperbolic manifolds},
Invent. Math. 84 (1986), no. 3, 523--540,
\href{https://doi.org/10.1007/BF01388745}{doi:10.1007/BF01388745}.

\bibitem[CDDP]{CDDP} M. \v{C}eki\'c, B. Delarue, S. Dyatlov, and G. P. Paternain,
\emph{The Ruelle zeta function at zero for nearly hyperbolic 3-manifolds},
Invent. Math. \textbf{229} (2022), no.~1, 303--394.
\href{https://doi.org/10.1007/s00222-022-01108-x}{doi:10.1007/s00222-022-01108-x};
\href{https://arxiv.org/abs/2009.08558v3}{arXiv:2009.08558v3}.
Equation references follow arXiv version~v3.

\bibitem[CD]{CD} Y. Chaubet and N. V. Dang,
\emph{Dynamical torsion for contact Anosov flows},
Anal. PDE \textbf{17} (2024), no.~8, 2619--2681.
\href{https://doi.org/10.2140/apde.2024.17.2619}{doi:10.2140/apde.2024.17.2619};
\href{https://arxiv.org/abs/1911.09931v2}{arXiv:1911.09931v2}.
\href{https://arxiv.org/html/1911.09931v2}{Author version~v2}
(equation numbering used here).

\bibitem[DGRS]{DGRS} N. V. Dang, C. Guillarmou, G. Rivi\`ere, and S. Shen,
\emph{The Fried conjecture in small dimensions},
Invent. Math. \textbf{220} (2020), no.~2, 525--579.
\href{https://doi.org/10.1007/s00222-019-00935-9}{doi:10.1007/s00222-019-00935-9};
\href{https://arxiv.org/abs/1807.01189v3}{arXiv:1807.01189v3}.

\bibitem[DR]{DR} N. V. Dang, G. Rivi\`ere,
\emph{Topology of Pollicott--Ruelle resonant states},
Ann. Sc. Norm. Super. Pisa Cl. Sci. (5) 21 (2020), special issue,
827--871,
\href{https://doi.org/10.2422/2036-2145.201804_010}{doi:10.2422/2036-2145.201804\_010},
\href{https://arxiv.org/abs/1703.08037}{Primary preprint}, Theorem~2.1.

\bibitem[MS]{MS} D. Margalit and S. Spallone,
\emph{A homological recipe for pseudo-Anosovs},
Math. Res. Lett. \textbf{14} (2007), no.~5, 853--863,
\href{https://doi.org/10.4310/MRL.2007.v14.n5.a12}{DOI};
\href{https://margalit.droppages.net/papers/homcrit.pdf}{author PDF},
Proposition~1.

\bibitem[T]{Thurston} W. P. Thurston,
\emph{Hyperbolic structures on 3-manifolds, II: surface groups and
3-manifolds which fiber over the circle},
\href{https://arxiv.org/abs/math/9801045}{Primary preprint}.

\bibitem[BO]{BO} U. Bunke and M. Olbrich,
\emph{Theta and zeta functions for odd-dimensional locally symmetric
spaces of rank one}, preprint, 1994,
\href{https://arxiv.org/abs/dg-ga/9407012}{arXiv:dg-ga/9407012},
\href{https://arxiv.org/html/dg-ga/9407012v1}{version 1},
Proposition 4.7 and its graded-determinant convention.
Related published treatment: \emph{Selberg Zeta and Theta Functions:
A Differential Operator Approach}, Mathematical Research \textbf{83},
Akademie Verlag, Berlin, 1995.
\end{thebibliography}
\end{document}
